
\section{The reductive center}
\SGARefTarget{sga3:xii:section:4}

\paragraph{\textbf{Definition 4.1.}}
\SGARefTarget{sga3:xii:4.1}
Let \(S\) be a prescheme, let \(G\) be an \(S\)-group of finite
presentation over \(S\) with affine fibers, and let \(Z\) be a group
subprescheme of \(G\). We say that \(Z\) is a \emph{reductive center}
of \(G\) if:

\begin{enumerate}
\item[(i)] \(Z\) is central and of multiplicative type;
\item[(ii)] for every base change \(S'\to S\) and every central
homomorphism
\[
  u:H\longrightarrow G_{S'},
\]
where \(H\) is a group of multiplicative type and of finite type over
\(S'\), the homomorphism \(u\) factors through \(Z_{S'}\).
\end{enumerate}

(For a variant of this notion when the fibers of \(G\) are not
assumed affine, see \SGARefLink{sga3:xii:8.6}{8.6}.)

Such a \(Z\) is necessarily of finite type over \(S\), since its
fibers are, and is therefore uniquely determined as the largest
central subgroup of multiplicative type in \(G\). It is easy to give
examples in which \(G\), smooth and affine over \(S\), admits a
largest central subgroup \(Z\) of multiplicative type but \(Z\) is
not a reductive center---that is, \(G\) admits no reductive center;
see, for example, \SGARefLink{sga3:xii:1.6}{1.6}. It follows readily, however, from \SGARefLink{sga3:ix:theorem:6:8}{IX~6.8} that
a subgroup \(Z\) of \(G\) is a reductive center if and only if it is
a largest central subgroup of multiplicative type and retains this
property after every base change. See also \SGARefLink{sga3:xii:4.3}{4.3} below.

It is clear from \SGARefLink{sga3:xii:4.1}{4.1} that if \(Z\) is the reductive center of \(G\),
then \(Z_{S'}\) is the reductive center of \(G_{S'}\) after every base
change \(S'\to S\). This fact and the uniqueness of the reductive
center, together with faithfully flat quasi-compact descent
(SGA~1, VIII), give:

\paragraph{\textbf{Proposition 4.2.}}
\SGARefTarget{sga3:xii:4.2}
\emph{Let \(G\) be an \(S\)-group prescheme of finite presentation
over \(S\), with affine fibers. If \(G\) admits a reductive center
locally for the faithfully flat quasi-compact topology, then it
admits a reductive center. For \(Z\) to be a reductive center of
\(G\), it is necessary and sufficient that it be so locally for the
fpqc topology.}

We also record:

\paragraph{\textbf{Proposition 4.3.}}
\SGARefTarget{sga3:xii:4.3}
\emph{Let \(G\) be an \(S\)-group prescheme of finite presentation
with affine fibers, and let \(Z\) be a group subprescheme. For \(Z\)
to be a reductive center of \(G\), it is necessary and sufficient
that it be of multiplicative type and that, for every \(s\in S\),
\(Z_s\) be a reductive center of \(G_s\).}

Indeed, \SGARefLink{sga3:ix:theorem:5:6:item:b}{IX~5.6(b)} first shows that \(Z\) is central. Since the
conditions under consideration are stable under base change, it
remains to prove that every central homomorphism
\[
  u:H\longrightarrow G,
\]
with \(H\) of multiplicative type and of finite type, factors through
\(Z\). Since \(Z\) is central, \(u\) and the canonical immersion
\(v:Z\to G\) define a group homomorphism
\[
  w:H\times_S Z\longrightarrow G.
\]
By \SGARefLink{sga3:ix:theorem:6:8}{IX~6.8}, this homomorphism has a group image \(K\), which is a
subgroup of multiplicative type in \(G\). Everything reduces to
showing that \(K=Z\). This holds fiber by fiber by the hypothesis on
\(Z\), and \SGARefLink{sga3:ix:theorem:5:1:bis}{IX~5.1 bis} now applies, completing the proof of \SGARefLink{sga3:xii:4.3}{4.3}.

In criterion \SGARefLink{sga3:xii:4.3}{4.3}, the hypothesis that \(Z_s\) be the reductive center
of \(G_s\) for every \(s\in S\) is in fact purely geometric: it is
enough to verify it after passing to an algebraic closure of
\(\kappa(s)\), as follows from the second assertion of \SGARefLink{sga3:xii:4.2}{4.2}.

\paragraph{\textbf{Theorem 4.4.}}
\SGARefTarget{sga3:xii:4.4}
\emph{Let \(G\) be an affine algebraic group over a field \(k\).
Then \(G\) admits a reductive center.}

Since the center of \(G\) is represented by a closed subgroup of
\(G\) by \SGARefLink{sga3:viii:corollary:6:7}{VIII~6.7}, we are immediately reduced to the case in which
\(G\) is commutative. Moreover, by \SGARefLink{sga3:xii:4.2}{4.2} we may suppose that \(k\) is
algebraically closed. We shall see in Exposé XVII that in this case
\(G\) can be written as a product \(Z\times U\), where \(Z\) is of
multiplicative type and \(U\) is ``unipotent'', from which it follows
at once that \(Z\) is a reductive center of \(G\). We give here a
proof independent of the general structure theorem for commutative
algebraic groups just indicated.

By \SGARefLink{sga3:ix:theorem:6:8}{IX~6.8}, the set of subgroups \(H\) of multiplicative type in
\(G\) is filtered upward. We show that it has a largest element.

When \(G\) is smooth over \(k\), the structure theorem
BIBLE~4, Th.~4 gives
\[
  G=G_s\times G_u,
\]
where \(G_s\) is of multiplicative type and \(G_u\) is ``unipotent'':
here this means that \(G_u\) has a composition series whose factors
are subgroups of \(\mathbb G_a\), smooth if desired. Indeed,
\(G_u/G_u^0\) is a unipotent group by BIBLE~4, Corollary to Th.~3,
hence a \(p\)-group, where \(p\) is the characteristic exponent, by
BIBLE~4, Prop.~4. On the other hand, \(G_u^0\) has a composition
series with smooth connected one-dimensional quotients by BIBLE~6,
Th.~1, Cor.~1, and these quotients are isomorphic to \(\mathbb G_a\)
by BIBLE~7, Th.~4. Every homomorphism from a group of multiplicative
type to \(\mathbb G_a\), and therefore also to \(G_u\), is trivial,
as follows from the lemma below. Thus every subgroup of
multiplicative type in \(G\) is contained in \(G_s\).

In the general case, consider the subgroup
\[
  G_0=G_{\mathrm{red}}
\]
of \(G\). It is smooth over \(k\), and hence by the preceding case
has a largest subgroup \(Z_0\) of multiplicative type. The subgroups
of multiplicative type in \(G\) containing \(Z_0\) correspond to
the subgroups of multiplicative type in \(G'=G/Z_0\). We have an
exact sequence
\[
  0\longrightarrow G_0/Z_0
   \longrightarrow G/Z_0
   \longrightarrow G/G_0
   \longrightarrow 0,
\]
where \(G_0/Z_0\) has no subgroup of multiplicative type other than
the unit group. Since every subgroup of a group of multiplicative
type is again of multiplicative type by IX~8, every subgroup \(H\)
of multiplicative type in \(G/Z_0\) satisfies
\[
  H\cap(G_0/Z_0)=0;
\]
thus \(H\to G/G_0\) is injective. The algebraic group \(G/G_0\) is
radicial, hence finite over \(k\), so \(H\) is itself radicial and
its rank is bounded above by that of \(G/G_0\). Consequently the
family of subgroups of multiplicative type in \(G\) has a largest
element, say \(Z\).

We claim that \(G/Z\) has no subgroup of multiplicative type other
than the unit subgroup. Indeed, a commutative extension of two
algebraic groups of multiplicative type is of multiplicative type
(\SGARefLink{sga3:xvii:proposition:7:1:1}{7.1.1}). If \(Z'\) is the inverse image in \(G\) of a subgroup of
multiplicative type in \(G/Z\), then \(Z'\) is of multiplicative
type. The maximality of \(Z\) gives \(Z'=Z\), and hence \(Z'/Z=0\).

The ``principle of finite extension'' now readily shows that for
every extension \(K\) of \(k\),
\[
  (G/Z)_K=G_K/Z_K
\]
likewise has no subgroup of multiplicative type other than the unit
group.

We can now prove that \(Z\) is a reductive center of \(G\). Let
\[
  u:H\longrightarrow G_S
\]
be a homomorphism of \(S\)-groups, where \(S\) is a prescheme over
\(k\) and \(H\) is of multiplicative type and of finite type over
\(S\). We must show that \(u\) factors through \(Z_S\), or
equivalently that the composite
\[
  H\longrightarrow G_S\longrightarrow
  (G/Z)_S=G_S/Z_S
\]
is trivial. Put \(U=G/Z\). Every homomorphism
\(H\to U_S\) is trivial: by \SGARefLink{sga3:ix:corollary:5:2}{IX~5.2} one reduces to the case where
\(S\) is the spectrum of a field \(K\), and \SGARefLink{sga3:ix:theorem:6:8}{IX~6.8} then reduces the
claim to the fact that \(U_K\) has no nontrivial subgroup of
multiplicative type. This proves \SGARefLink{sga3:xii:4.4}{4.4}, modulo the following lemma.

\paragraph{\textbf{Lemma 4.4.1.}}
\SGARefTarget{sga3:xii:4.4.1}
\emph{Let \(H\) be an \(S\)-group prescheme of multiplicative type.
Then every homomorphism of \(S\)-groups
\[
  u:H\longrightarrow\mathbb G_a
\]
is trivial.}

Regard \(E=\mathcal O_S^2\) as an extension with submodule
\(E'\simeq\mathcal O_S\) and quotient \(E''\simeq\mathcal O_S\).
Then \(\mathbb G_a\) identifies with the prescheme of automorphisms
of this extension. A homomorphism \(u:H\to\mathbb G_a\) therefore
amounts to an \(H\)-module structure on \(E\) preserving the
extension, with \(E'\) stable under \(H\) and with the induced
actions on \(E'\) and \(E''\) trivial. By \SGARefLink{sga3:i:proposition:4-7-3}{I~4.7.3}, \(H\) acts
trivially on \(E\); hence \(u\) is trivial.

\paragraph{\textbf{Remark 4.4.2.}}
\SGARefTarget{sga3:xii:4.4.2}
If \(G\) is a nonzero abelian variety over \(k\), then \(G\) admits
no reductive center in the sense of \SGARefLink{sga3:xii:4.1}{4.1} with the restriction
``\(G\) affine'' omitted. Indeed, if \(n\) is prime to the
characteristic, then \({}_nG\) is étale over \(k\), of order prime
to the characteristic, and hence of multiplicative type. The family
of the \({}_nG\) is schematically dense in \(G\). A reductive center,
if it existed, would therefore equal \(G\), which is impossible
because \(G\) is not of multiplicative type. This explains the
requirement in \SGARefLink{sga3:xii:4.1}{4.1} that \(G\) have affine fibers.

\paragraph{\textbf{Lemma 4.5.}}
\SGARefTarget{sga3:xii:4.5}
\emph{Let \(S\) be a prescheme, let \(G\) be a smooth affine
\(S\)-group prescheme with connected fibers, let \(T\) be a maximal
torus of \(G\), and let \(u:H\to G\) be a central homomorphism, where
\(H\) is an \(S\)-group prescheme of multiplicative type and of
finite type. Then \(u\) factors through \(T\).}

Let \(C=\operatorname{Centr}_G(T)\). By \SGARefLink{sga3:xi:corollary:5:3}{XI~5.3} it is a closed smooth
group subprescheme of \(G\), hence affine over \(S\). Since \(T\) is
central in \(C\), it is invariant, and the quotient group
\[
  U=C/T
\]
is represented by \SGARefLink{sga3:viii:theorem:5:1}{VIII~5.1}. Because \(u\) is central it factors
through \(C\), and it remains only to show that the composite
\[
  H\longrightarrow C\longrightarrow U=C/T
\]
is trivial. By \SGARefLink{sga3:ix:corollary:5:2}{IX~5.2} we reduce to the case where \(S\) is the
spectrum of a field, which we may suppose algebraically closed.
Then BIBLE~6, Th.~2 shows that \(U\) is a connected algebraic group,
smooth over \(k\), which is ``unipotent''; as recalled above, it has
a composition series with quotients isomorphic to \(\mathbb G_a\).
Every homomorphism from the finite-type group \(H\) of multiplicative
type to \(U\) is therefore trivial. This proves \SGARefLink{sga3:xii:4.5}{4.5}.

\paragraph{\textbf{Corollary 4.6.}}
\SGARefTarget{sga3:xii:4.6}
\emph{Let \(G\) be as in \SGARefLink{sga3:xii:4.5}{4.5}. If \(G\) admits a reductive center,
then that center is contained in every maximal torus of \(G\).}

\paragraph{\textbf{Theorem 4.7.}}
\SGARefTarget{sga3:xii:4.7}
\emph{Let \(S\) be a prescheme and let \(G\) be an \(S\)-group
prescheme, smooth and affine over \(S\), with connected fibers.}

\begin{enumerate}
\item[\textup{\SGAInlineRefTarget{sga3:xii:4.7:item:a}(a)}] \emph{For every \(s\in S\), let \(z(s)\) denote
the type of the reductive center of \(G_s\), which is defined by
\SGARefLink{sga3:xii:4.4}{4.4}. Order the set \(E\) of isomorphism classes of finite-type
\(\mathbb Z\)-modules by declaring the class of \(M\) to be greater
than that of \(M'\) when \(M'\) is isomorphic to a quotient of \(M\).
Then the map
\[
  S\longrightarrow E,\qquad s\longmapsto z(s),
\]
is lower semicontinuous.}

\item[\textup{\SGAInlineRefTarget{sga3:xii:4.7:item:b}\SGAInlineRefTarget{sga3:xii:theorem:4-7:item:b}(b)}] \emph{The group \(G\) admits a reductive center
\(Z\) if and only if the preceding function \(z:S\to E\) is locally
constant. If this holds, then \(G/Z\) is representable
\textup{(cf.\ \SGARefLink{sga3:viii:theorem:5:1}{VIII~5.1})}, and its reductive center is the unit
subgroup.}

\item[\textup{\SGAInlineRefTarget{sga3:xii:4.7:item:c}\SGAInlineRefTarget{sga3:xii:theorem:4-7:item:c}(c)}] \emph{Suppose that the reductive rank of \(G\)
is locally constant \textup{(cf.\ \SGARefLink{sga3:xii:1.7:item:b}{1.7(b)})}. Then \(G\) admits a
reductive center \(Z\). If \(G/Z\) is representable---for example,
if \(G\) is affine over \(S\)---then the maximal tori \(T\) of \(G\)
\textup{(respectively, the Cartan subgroups \(C\) of \(G\))} and the
maximal tori \(T'\) of \(G'=G/Z\) \textup{(respectively, the Cartan
subgroups \(C'\) of \(G'\))} are in bijective correspondence:
\[
  T'=T/Z,\qquad C'=C/Z,
\]
and conversely
\[
  T=\varphi^{-1}(T'),\qquad C=\varphi^{-1}(C'),
\]
where \(\varphi:G\to G'\) is the canonical homomorphism.}

\item[\textup{\SGAInlineRefTarget{sga3:xii:4.7:item:d}(d)}] \emph{Let \(T\) be a maximal torus of \(G\), let
\(\mathfrak g\) denote the Lie algebra of \(G\), and consider the
homomorphism
\[
  \theta:T\longrightarrow\operatorname{GL}(\mathfrak g)
\]
induced by the adjoint representation of \(G\) \textup{(II~4)}.
Then \(\ker\theta\) is a reductive center of \(G\).}
\end{enumerate}

\emph{Proof.}
\textup{\SGARefLink{sga3:xii:4.7:item:a}{(a)}} and \textup{\SGARefLink{sga3:xii:4.7:item:b}{(b)}}. The proof of \textup{\SGARefLink{sga3:xii:4.7:item:a}{(a)}} and of the
first assertion of \textup{\SGARefLink{sga3:xii:4.7:item:b}{(b)}} is essentially identical to that of
\SGARefLink{sga3:xii:1.7:item:a}{1.7(a)} and \SGARefLink{sga3:xii:4.7:item:b}{(b)}, and we shall not repeat the argument. We point out
only that in \textup{\SGARefLink{sga3:xii:4.7:item:a}{(a)}} one must appeal to \SGARefLink{sga3:ix:theorem:5:6}{IX~5.6}, using the fact
that \(G\) has connected fibers.

We prove the second assertion of \textup{\SGARefLink{sga3:ix:theorem:5:6:item:b}{(b)}}. Suppose that \(Z\) is
a reductive center of \(G\), and put \(G'=G/Z\). By \SGARefLink{sga3:viii:theorem:5:1}{VIII~5.1} the
quotient is representable. It is affine over \(S\), of finite
presentation over \(S\) by \SGARefLink{sga3:viii:corollary:5:8}{VIII~5.8}, and even smooth over \(S\):
this may be checked on the fibers because \(G'\) is flat and of
finite presentation over \(S\), and VI\textsubscript{B}~\SGARefLink{concl:sga3-VIB-9.2-xii}{9.2(xii)}
states that a quotient of a smooth algebraic group over a field is
smooth. Since \(G\) has connected fibers, so does \(G'\). Thus \(G'\)
satisfies the same hypotheses as \(G\).

By \SGARefLink{sga3:xii:4.3}{4.3} we may reduce to the case in which \(S\) is the spectrum of a
field. Let \(Z'\) be a central subgroup of multiplicative type in
\(G'\); it remains to prove that \(Z'\) is the unit subgroup. Let
\(Z_1\) be the inverse image of \(Z'\) in \(G\), and consider the
action of \(Z_1\) on \(G\) by inner automorphisms. Since \(Z\) is
central, it acts trivially, and hence \(Z_1\) acts through an action
of \(Z'\) on \(G\). Moreover \(Z'\) acts trivially both on \(Z\) and
on \(G/Z=G'\). Since \(Z\) is central, the action of \(Z'\) on \(G\)
therefore comes from a group homomorphism
\[
  u:Z'\longrightarrow
    \operatorname{Hom}_{S\text{-}\mathrm{gr}}(G',Z),
\]
characterized by
\[
  r(z')\cdot g
    =g\cdot u(z')\bigl(\varphi(g)\bigr),
\]
where \(r:Z'\to\operatorname{Aut}_{S\text{-}\mathrm{gr}}(G)\) is
the representation just considered and \(\varphi:G\to G'\) is the
canonical homomorphism. Giving \(u\) is equivalent to giving a group
homomorphism
\[
  v:G'\longrightarrow
    \operatorname{Hom}_{S\text{-}\mathrm{gr}}(Z',Z).
\]
By \SGARefLink{sga3:x:corollary:5-8}{X~5.8} the group on the right is representable and étale over
\(S\). Its unit section is therefore an open immersion, so
\(\ker v\) is an open subgroup of \(G'\). Since \(G'\) has connected
fibers, \(\ker v=G'\). Thus \(v=0\), hence \(u=0\), and \(Z'\), and
therefore \(Z_1\), acts trivially on \(G\). Consequently \(Z_1\) is
central in \(G\). It is a commutative extension of the group \(Z'\)
of multiplicative type by the group \(Z\) of multiplicative type,
and hence, over a field, is itself of multiplicative type by \SGARefLink{sga3:xvii:proposition:7:1:1}{7.1.1}.
Being central in \(G\), it is contained in the reductive center
\(Z\). Thus \(Z_1=Z\), whence \(Z'\) is the unit subgroup. This
proves \textup{\SGARefLink{sga3:xii:4.7:item:b}{(b)}}.

\textup{\SGARefLink{sga3:xii:4.7:item:d}{(d)}} Put \(Z=\ker\theta\), a subgroup of multiplicative type
in \(T\), for example by \SGARefLink{sga3:ix:theorem:6:8}{IX~6.8}. By \SGARefLink{sga3:xii:4.5}{4.5}, every central homomorphism
\(u:H\to G\), with \(H\) of multiplicative type and of finite type,
factors through \(T\), and therefore through \(Z\). Since formation
of \(Z\) commutes with every base change, it remains to show that
\(Z\) is central, or equivalently that its centralizer \(C\) is
equal to \(G\). By \SGARefLink{sga3:xi:corollary:5:3}{XI~5.3}, \(C\) is a closed smooth subgroup of
\(G\). Since \(Z\) acts trivially on \(\mathfrak g\), we have
\[
  \operatorname{Lie}(C)=\mathfrak g.
\]
It follows that \(C=G\). Indeed, the immersion \(C\to G\) is étale,
as may be checked fiberwise: it is an unramified homomorphism of
smooth algebraic groups of the same dimension
\textup{(VI\textsubscript{B}~\SGARefLink{sga3:xii:1.3}{1.3})}, so \SGARefLink{sga3:x:lemma:3:5}{X~3.5} applies. Hence
\(C\to G\) is an étale closed immersion, and therefore an open
immersion by SGA~1, \SGARefLink{sga3:i:subsection:5-1}{I~5.1}. Since it is both open and closed and \(G\)
has connected fibers, it is an isomorphism.

\textup{\SGARefLink{sga3:xii:4.7:item:c}{(c)}} By 1.7(b), \(G\) admits a maximal torus locally for the
fpqc topology. \SGARefLink{sga3:xii:4.2}{Proposition~4.2} and part \textup{\SGARefLink{sga3:xii:corollary:6-7:equation:d}{(d)}} therefore show
that \(G\) admits a reductive center \(Z\). Part \textup{\SGARefLink{sga3:xii:corollary:6-7:equation:d}{(d)}} also
shows that every maximal torus \(T\) of \(G\) contains \(Z\).
Clearly \(T'=T/Z\) is a torus. To prove that it is a maximal torus
of \(G'\), \SGARefLink{sga3:xii:1.3}{Definition~1.3} reduces us to the case in which \(S\) is
the spectrum of an algebraically closed field; the assertion then
follows from \emph{BIBLE}~7, Th.~3\SGARefLink{sga3:xii:4.7:item:a}{(a)}.

Conversely, let \(T'\) be a maximal torus of \(G'\), and set
\(T=\varphi^{-1}(T')\). We prove that \(T\) is a maximal torus of
\(G\). The question is local for the fpqc topology, so we may
suppose that \(G\) already admits a maximal torus \(T_0\). By what
precedes, \(T'_0=T_0/Z\) is a maximal torus of \(G'\). By the
conjugacy theorem 1.7(c), \(T'\) and \(T'_0\) are locally conjugate
for the fpqc topology. Since \(\varphi:G\to G'\) is a covering for
this topology, sections of \(G'\) lift locally to sections of \(G\);
thus \(T\) and \(T_0\) are likewise locally conjugate. Hence \(T\)
is a maximal torus. The argument for Cartan subgroups is identical.
This completes the proof.

Here is a useful reformulation of \textup{\SGARefLink{sga3:xii:corollary:6-7:equation:d}{(d)}} when \(T\) is
diagonalizable, say \(T=D_S(M)\), where \(M\) is a finite-type free
\(\mathbb Z\)-module. By \SGARefLink{sga3:i:proposition:4-7-3}{I~4.7.3}, the \(T\)-module \(\mathfrak g\)
decomposes into the direct sum of its \(T\)-submodules
\(\mathfrak g_m\), \(m\in M\):
\[
  \mathfrak g=\bigoplus_{m\in M}\mathfrak g_m.
\]
These submodules are necessarily locally free. Suppose that for
every \(m\in M\) the rank of \(\mathfrak g_m\) is constant, as is
the case in particular when \(S\) is connected. Then the set
\[
  R=\{m\in M\mid\mathfrak g_m\ne0\}
\]
of ``roots'' is finite. We obtain:

\paragraph{\textbf{Corollary 4.8.}}
\SGARefTarget{sga3:xii:4.8}
\emph{Under the preceding hypotheses and with the preceding
notation, the reductive center of \(G\) is the intersection of the
kernels in \(T\) of the root characters \(m\in R\). Thus there is
an isomorphism
\[
  Z\simeq D_S(N),
\]
where \(N\) is the quotient of \(M\) by the subgroup generated by
\(R\).}

\paragraph{\textbf{Corollary 4.9.}}
\SGARefTarget{sga3:xii:4.9}
\emph{Let \(G\) be an algebraic group over an algebraically closed
field \(k\). Suppose that \(G\) is smooth, connected, and affine
over \(k\), that its reductive center is the unit subgroup, and that
its Lie algebra \(\mathfrak g\) is nilpotent. Then \(G\) is
``unipotent,'' that is, it admits a composition series whose
quotients are isomorphic to \(\mathbb G_a\).}

By \emph{BIBLE}~6, Th.~4, Cor.~3, it is enough to prove that a
maximal torus \(T\) of \(G\) is the unit subgroup, or equivalently
that its Lie algebra \(\mathfrak t\) is zero. The \(T\)-module
\(\mathfrak g\) decomposes according to the characters of \(T\) by
\SGARefLink{sga3:i:proposition:4-7-3}{I~4.7.3}, so for every \(t\in\mathfrak t\), the endomorphism
\(\operatorname{ad}_{\mathfrak g}(t)\) of \(\mathfrak g\) is
semisimple. If \(\mathfrak g\) is nilpotent, this endomorphism is
also nilpotent, hence zero. By \SGARefLink{sga3:xii:4.7:item:d}{4.7(d)}, the hypothesis on the
reductive center makes the homomorphism
\[
  T\longrightarrow\operatorname{GL}(\mathfrak g)
\]
a monomorphism, and it therefore induces an injective map on Lie
algebras. By \SGARefLink{sga3:ii:proposition:4-5}{II~4.5}, this says that
\(\operatorname{ad}_{\mathfrak g}(t)=0\) implies \(t=0\).
Consequently \(\mathfrak t=0\), as required.

\paragraph{\textbf{Proposition 4.10.}}
\SGARefTarget{sga3:xii:4.10}
\emph{Let \(G\) be a smooth connected affine algebraic group over
an algebraically closed field \(k\). Then the reductive center of
\(G\) is the intersection of the maximal tori of \(G\).}

The intersection is meant schematically, or equivalently in the
sense of subfunctors of \(G\): it is the largest closed
subprescheme of \(G\) contained in every maximal torus. Since \(G\)
is noetherian, it is already the intersection of a suitable finite
family of maximal tori.

Let \(Z\) denote this intersection. It is a closed subgroup of a
torus, hence is of multiplicative type, and by \SGARefLink{sga3:xii:4.5}{4.5} it contains the
reductive center of \(G\). To prove equality it remains to show that
\(Z\) is central. Since \(G\) is connected, \SGARefLink{sga3:ix:corollary:5:5}{IX~5.5} reduces this to
showing that \(Z\) is invariant. By construction \(Z\) is invariant
under every inner automorphism \(\operatorname{int}(g)\), with
\(g\in G(k)\). Its normalizer \(N\), represented by \SGARefLink{sga3:viii:corollary:6:7}{VIII~6.7}, is
therefore a closed subgroup of \(G\) containing all rational points
of \(G\). Since \(G\) is reduced, \(N=G\), which proves the result.

\paragraph{\textbf{Proposition 4.11.}}
\SGARefTarget{sga3:xii:4.11}
\emph{Let \(S\) be a prescheme and let \(G\) be a smooth affine
\(S\)-group prescheme with connected fibers and unipotent rank zero.
Then \(G\) has locally constant reductive rank by \SGARefLink{sga3:xii:1.9}{1.9}, and hence the
``prescheme of maximal tori of \(G\),'' denoted by
\(\mathcal T\), is defined by \SGARefLink{sga3:xii:1.10}{1.10} and is smooth, separated, and of
finite type over \(S\). Let \(G\) act on \(\mathcal T\) through
inner automorphisms, giving a homomorphism of group functors
\[
  u:G\longrightarrow\operatorname{Aut}_S(\mathcal T).
\]
Under these hypotheses the following three subfunctors of \(G\) are
identical:
\begin{enumerate}
\item[\textup{(i)}] the reductive center \(Z\) of \(G\);
\item[\textup{(ii)}] the center \(Z'\) of \(G\);
\item[\textup{(iii)}] the kernel \(Z''=\ker u\) of the preceding
homomorphism.
\end{enumerate}
In particular, the center of \(G\) is represented by a subgroup of
multiplicative type in \(G\).}

\emph{Proof.}
Clearly \(Z\subset Z'\subset Z''\). It remains to prove
\(Z''\subset Z\), or, since the hypotheses are stable under base
change, that every section \(g\) of \(G\) over \(S\) which acts
trivially on \(\mathcal T\) is a section of \(Z\). Put \(G'=G/Z\).
By \SGARefLink{sga3:xii:4.7:item:b}{4.7(b)} and \SGARefLink{sga3:xii:4.7:item:c}{(c)}, the prescheme \(\mathcal T'\) of maximal tori of
\(G'\) is canonically isomorphic to \(\mathcal T\). We are therefore
reduced to the case \(G=G'\), in which the reductive center \(Z\) is
the unit subgroup. Notice that \SGARefLink{sga3:xii:4.7:item:c}{4.7(c)} also shows that the unipotent
rank of \(G'\) equals that of \(G\), and hence is zero.

We must prove that \(g\) is the unit section of \(G\). The usual
reduction process reduces us first to the case in which \(S\) is
noetherian, and then to the case in which \(S\) is local artinian:
when \(S\) is locally noetherian, it is enough to check the assertion
after every base change \(S'\to S\) with \(S'\) local artinian. In
this case \(Z''=\ker u\) is representable by \SGARefLink{sga3:viii:proposition:6:2:item:a:01}{VIII~6.2(a)} and \SGARefLink{sga3:viii:examples:6:5:item:c:01}{6.5(c)};
alternatively, representability follows from \SGARefLink{sga3:xi:theorem:6:8}{XI~6.8}. By Nakayama's
lemma it is enough to show that the fiber \(Z''_0\) is the unit
subgroup. We may therefore suppose that \(S\) is the spectrum of a
field \(k\), which we may take to be algebraically closed.

The subgroup \(Z''\) lies in the stabilizer of every point of
\(\mathcal T(k)\), that is, in the normalizer \(N(T)\) of every
maximal torus \(T\) of \(G\). Since the unipotent rank of \(G\) is
zero, \SGARefLink{sga3:xi:proposition:5:9}{XI~5.9} says that \(T\) is an open subgroup of \(N(T)\). Thus
\[
  \operatorname{Lie}N(T)=\operatorname{Lie}T,
\]
and \(\operatorname{Lie}Z''\) is contained in
\(\operatorname{Lie}T\) for every maximal torus \(T\). \SGARefLink{sga3:xii:4.10}{Proposition
4.10} shows that the intersection of the Lie algebras of all maximal
tori is the Lie algebra of the reductive center \(Z\), which is zero
in the present reduction. Hence \(\operatorname{Lie}Z''=0\), so
\(Z''\) is étale over \(k\).

Moreover, \(Z''\) is invariant in \(G\). Since \(G\) is connected,
it follows that \(Z''\) is central in \(G\). Thus it is contained in
\[
  T=\operatorname{Centr}_G(T)
\]
for every maximal torus \(T\), and therefore in their intersection.
By \SGARefLink{sga3:xii:4.10}{4.10} this intersection is \(Z=0\), completing the proof.

\section[Application to the prescheme of subgroups of multiplicative type]
{Application to the prescheme of subgroups of multiplicative
type\texorpdfstring{\protect\footnotemark}{}}
\footnotetext{For a generalization of the results of this section to
the case in which \(G\) is not assumed affine over \(S\), see
M.~Raynaud, ``Ample sheaves on group schemes and homogeneous
spaces,'' Chapter~IX.}
\SGARefTarget{sga3:xii:5}

\paragraph{\textbf{Theorem 5.1.}}
\SGARefTarget{sga3:xii:5.1}
\emph{Let \(S\) be a prescheme, let \(G\) be a smooth affine
\(S\)-group prescheme, and let \(M\) be the ``prescheme of subgroups
of multiplicative type in \(G\),'' representing the functor
described in \SGARefLink{sga3:xi:theorem:4:1}{XI~4.1}. For every integer \(n>0\), let \(T_n\) be the
subfunctor of \(M\) for which \(T_n(S')\) is the set of subgroup
preschemes \(H\) of multiplicative type in \(G_{S'}\) satisfying
\[
  n\operatorname{id}_H=0.
\]
Thus \(T_n\) is representable and affine over \(S\) by \SGARefLink{sga3:xi:proposition:3:12:item:a}{XI~3.12(a)}.
Let
\[
  u_n:M\longrightarrow T_n
\]
be the morphism defined by \(u_n(H)={}_nH\), where
\({}_nH=\ker(n\operatorname{id}_H)\). Then:}

\begin{enumerate}
\item[\textup{\SGAInlineRefTarget{sga3:xii:theorem:5-1:item:a}\SGAInlineRefTarget{sga3:xii:5.1:item:a}(a)}] \emph{Every subprescheme \(U\) of \(M\) of
finite type over \(S\) is contained in a closed subprescheme of
finite type over \(S\). Every closed subprescheme \(X\) of \(M\) of
finite type over \(S\) is affine over \(S\).}

\item[\textup{\SGAInlineRefTarget{sga3:xii:5.1:item:b}(b)}] \emph{Suppose that \(S\) is quasi-compact, and
let \(X\) be a closed subprescheme of \(M\) of finite type over
\(S\). Then there is an integer \(n>0\) such that, for every
multiple \(m\) of \(n\), the induced morphism
\[
  u_m|_X:X\longrightarrow T_m
\]
is a closed immersion.}
\end{enumerate}

\emph{Proof.}
\textup{(a)} For the first assertion, take for \(X\) the schematic
closure of \(U\) \textup{(EGA~I, 9.5.1 and 9.5.3)}. This closure is
defined because the immersion \(i:U\to M\) is quasi-compact:
\(U\) is of finite type over \(S\), and \(M\) is separated over
\(S\) by \SGARefLink{sga3:xi:theorem:4:1}{XI~4.1}. It remains to prove that such an \(X\) is affine
over \(S\), which will prove the second assertion at the same time.
The question is local on \(S\), so we may suppose that \(S\) is
affine. Since \(U\) is of finite type over \(S\), it is contained
in a quasi-compact open of \(M\). We are therefore reduced to the
case in which \(U\) itself is an open subprescheme of finite type
over \(S\), hence quasi-compact. Notice that a closed subscheme of a
scheme affine over \(S\) is affine over \(S\).

It is enough to show that such a \(U\) is contained in a closed
affine subprescheme of \(M\). In this form the usual reduction
process reduces us immediately to the case in which \(S\) is
noetherian. The same argument applies to \textup{\SGARefLink{sga3:xii:5.1:item:b}{(b)}}, reducing it
to the case in which \(S\) is affine noetherian.

Recall the inverse limit \(T\) of the \(T_n\) used in the proof of
\SGARefLink{sga3:xi:theorem:4:1}{XI~4.1}. It is an affine prescheme over \(S\), though not of finite
type, and its local rings are noetherian, as was observed at the
beginning of the proof of IX~3.7. More explicitly, if
\(t=(t_n)\in T\), then the transition morphisms \(T_m\to T_n\) are
smooth and the dimensions of the \(T_m\) are bounded. Hence there
is an \(n_0\) such that \(T_n\to T_{n_0}\) is étale at \(t_n\) for
every multiple \(n\) of \(n_0\). The proof just cited, or \SGARefLink{sga3:xi:corollary:3:11}{XI~3.11},
shows that the canonical morphism
\[
  u:M\longrightarrow T
\]
is an immersion and induces isomorphisms on local rings. It need not
be an open immersion or a quasi-compact morphism.

Let \(U\) be a quasi-compact open of \(M\). We shall show that its
closure in \(T\) is contained in \(M\). If \(X\) is the schematic
closure of \(U\) in \(M\), this will make
\[
  u|_X:X\longrightarrow T
\]
a closed immersion. Then \(X\) is affine, proving \textup{\SGARefLink{sga3:xii:5.1:item:a}{(a)}}.
Since \(U\) is noetherian it has only finitely many irreducible
components. Every \(t\in T\) in the closure of \(U\) is a
specialization of a point \(x\in U\). The local ring of \(t\) in
\(T\) is noetherian, and so is the local ring \(A\) of \(t\) in
\(\overline{\{x\}}\), endowed with the induced reduced structure.

The canonical morphism from the noetherian integral local scheme
\[
  S'=\operatorname{Spec}(A)
\]
to \(T\) sends its closed point to \(t\) and its generic point to
\(x\in M\). We must show that \(t\in M\). Replacing \(A\), if
necessary, by the quotient by a minimal prime of a suitable complete
local \(A\)-algebra flat over \(A\) \textup{(EGA
0\textsubscript{III}, 10.3.1)}, we may suppose that \(A\) is
complete and has algebraically closed residue field. One could even
reduce further to a discrete valuation ring by EGA~II, 7.1.7.
Renaming \(S'\) as \(S\), we are reduced to the following lemma.

\paragraph{\textbf{Lemma 5.2.}}
\SGARefTarget{sga3:xii:5.2}
\emph{Let \(S\) be the spectrum of a complete noetherian integral
local ring with algebraically closed residue field. Let \(G\) be a
smooth affine \(S\)-group prescheme, let
\(\bigl(H(n)\bigr)_{n>0}\) be a system of subgroups of
multiplicative type in \(G\), indexed by the positive integers, and
let \(\eta\) be the generic point of \(S\). Suppose that:}

\begin{enumerate}
\item[\textup{(a)}] \emph{if \(m\) is a multiple of \(n\), then
\[
  H(n)={}_nH(m);
\]}
\item[\textup{\SGAInlineRefTarget{sga3:xii:5.2:item:b}(b)}] \emph{there is a subgroup \(H_\eta\) of
multiplicative type in \(G_\eta\) such that
\[
  H(n)_\eta={}_n(H_\eta)
\]
for every \(n>0\).}
\end{enumerate}

\emph{Then there is a subgroup \(H\) of multiplicative type in
\(G\) such that \(H(n)={}_nH\) for every \(n>0\).}

\emph{Proof.}
For every \(n\), let
\[
  C_n=\operatorname{Centr}_G\bigl(H(n)\bigr).
\]
This is a closed subgroup prescheme of \(G\), smooth over \(S\), by
\SGARefLink{sga3:xi:corollary:5:3}{XI~5.3}. The positive integers being ordered by divisibility, the
\(C_n\) form a decreasing family of closed subpreschemes of \(G\).
Since \(G\) is noetherian, this family is stationary. Let \(C\) be
the common value of the \(C_n\) for sufficiently divisible \(n\).
Then \(C\) satisfies the same conditions as \(G\), and all the
\(H(n)\) are in fact subgroups of \(C\). Replacing \(G\) by \(C\),
we may therefore suppose that all the \(H(n)\) are central subgroups
of \(G\).

Let \(s\) be the closed point of \(S\), and let \(Z_s\) be the
reductive center of \(G_s\), defined in \SGARefLink{sga3:xii:4.4}{4.4}. By the definition in
\SGARefLink{sga3:xii:4.1}{4.1}, \(Z_s\) contains every \(H(n)_s\). Since \(A\) is complete,
\(Z_s\) lifts to a subgroup prescheme \(Z\) of multiplicative type
in \(G\), by \SGARefLink{sga3:xi:theorem:5:8}{XI~5.8}. We claim that \(Z\) contains every \(H(n)\).
Indeed, \(H(n)\) is central and hence commutes with \(Z\) in \(G\).
The resulting homomorphism
\[
  Z\times H(n)\longrightarrow G
\]
has, by \SGARefLink{sga3:ix:theorem:6:8}{IX~6.8}, an image \(K\) which is a subgroup of multiplicative
type in \(G\) containing \(Z\). It remains only to prove \(K=Z\);
this follows from \(K_s=Z_s\) and \SGARefLink{sga3:ix:theorem:5:1}{IX~5.1}~bis.

We are thus reduced to proving the analogue of the lemma with \(G\)
replaced by a group \(Z\) of multiplicative type and of finite type
over \(S\), not necessarily smooth over \(S\). Since \(A\) is
complete with algebraically closed residue field, \(Z\) is
diagonalizable by \SGARefLink{sga3:x:corollary:3:3}{X~3.3} and \SGARefLink{sga3:xii:1.4}{1.4}, hence is of the form
\(D_S(M)\), where \(M\) is a finitely generated commutative group.
Every subgroup of multiplicative type \(H\) of \(Z\) is therefore
diagonalizable by \SGARefLink{sga3:ix:proposition:2:11:item:i:01}{IX~2.11(i)}, and so is of the form \(D_S(N)\),
where \(N\) is a quotient group of \(M\), by \SGARefLink{sga3:viii:corollary:3:2}{VIII~3.2}. Thus the
\(H(n)\) correspond to quotient groups \(M(n)\) of \(M\), and the
existence of a subgroup \(H\) such that \(H(n)={}_nH\) for every
\(n\) is equivalent to the existence of a quotient group \(N\) of
\(M\) such that
\[
  M(n)=N/nN
\]
for every \(n\). This follows at once from the existence of the
subgroup \(H_\eta\) of \(G_\eta\) such that
\(H(n)_\eta={}_n(H_\eta)\) for every \(n\). This proves \SGARefLink{sga3:xii:5.2}{Lemma~5.2}
and hence \SGARefLink{sga3:xii:theorem:5-1:item:a}{Theorem~5.1(a)}.

\textup{\SGARefLink{sga3:xii:5.2:item:b}{(b)}} We already know that
\[
  u|_X:X\longrightarrow T
\]
is a closed immersion. It follows readily that, for all sufficiently
large \(m\), the composite
\[
  u_m|_X:X\longrightarrow T_m
\]
is also a closed immersion. Since we shall not need this fact below,
we omit the details of the proof.

\paragraph{\textbf{Corollary 5.3.}}
\SGARefTarget{sga3:xii:5.3}
\emph{With the notation of \SGARefLink{sga3:xii:5.1}{5.1}, let \(U\) be a subprescheme of
\(M\), both open and closed, and of finite type over \(S\). Then
\(U\), with the structure induced by \(M\), is affine over \(S\).
If \(S\) is quasi-compact, there is an integer \(n>0\) such that,
for every multiple \(m\) of \(n\), the induced morphism
\[
  u_m|_U:U\longrightarrow T_m
\]
is both an open and a closed immersion; equivalently, it is an
isomorphism onto an open-and-closed subprescheme of \(T_m\), endowed
with the induced structure.}

\emph{Proof.}
The first assertion follows from \SGARefLink{sga3:xii:5.1:item:a}{5.1(a)}, and the second from \SGARefLink{sga3:xii:5.1:item:b}{5.1(b)},
since \(u_m:M\to T_m\) is smooth by \SGARefLink{sga3:xi:corollary:2:2}{XI~2.2}~bis and a smooth
immersion, that is, an étale immersion, is an open immersion.

\paragraph{\textbf{Corollary 5.4.}}
\SGARefTarget{sga3:xii:5.4}
\emph{Let \(S\) be a prescheme, let \(G\) be a smooth affine
\(S\)-group prescheme whose reductive rank is locally constant
\textup{(\SGARefLink{sga3:xii:1.7:item:b}{1.7(b)})}, and let \(\mathcal T\) be the ``prescheme of
maximal tori of \(G\)'' \textup{(\SGARefLink{sga3:xii:1.10}{1.10})}. Then \(\mathcal T\) is
smooth and affine over \(S\). If \(T\) is a maximal torus of \(G\)
and \(N(T)\) its normalizer, then \(G/N(T)\) \textup{(\SGARefLink{sga3:xi:corollary:5:3}{XI~5.3}~bis)}
is affine over \(S\). The same is true of \(G/C(T)\), where \(C(T)\)
is the centralizer of \(T\), provided that
\[
  W(T)=N(T)/C(T)
\]
is finite over \(S\); cf.\ \SGARefLink{sga3:xii:2.1:item:b}{2.1(b)}.}

\emph{Proof.}
The second assertion is contained in the first, since the conjugacy
theorem identifies \(G/N(T)\) with \(\mathcal T\), by \SGARefLink{sga3:xi:remark:5:5}{XI~5.5}~bis.
For the first assertion, observe that by construction
\(\mathcal T\) is isomorphic to an open-and-closed subprescheme of
\(M\). Indeed, one may suppose that the reductive rank of \(G\) is
constant and equal to \(r\); then \(\mathcal T\) is the subprescheme
of \(M\) corresponding to the subtori of relative dimension \(r\),
that is, the largest subprescheme of \(M\) over which the universal
subgroup of multiplicative type
\[
  H\subset G_M
\]
is a torus of relative dimension \(r\). \SGARefLink{sga3:xii:5.3}{Corollary~5.3} therefore
applies. Finally, \(G/C(T)\) is finite over
\[
  G/N(T)\simeq \bigl(G/C(T)\bigr)/W(T),
\]
and hence is affine because \(G/N(T)\) is affine.

\paragraph{\textbf{Corollary 5.5.}}
\SGARefTarget{sga3:xii:5.5}
\emph{Let \(G\) be a smooth affine algebraic group over a field
\(k\). Then the scheme \(M\) of subgroups of multiplicative type in
\(G\) is a disjoint union of affine \(k\)-schemes. For every
subgroup \(H\) of multiplicative type in \(G\), if \(C\) and \(N\)
denote respectively its centralizer and its normalizer, then the
quotients \(G/C\) and \(G/N\) are affine.}

\emph{Proof.}
By \SGARefLink{sga3:xi:corollary:5:1}{XI~5.1}~bis, the saturation under the action of \(G\) on \(M\) of
every finite closed subset of \(M\) is open. Indeed, one reduces to
the case in which \(k\) is algebraically closed, and hence to the
orbit of a \(k\)-rational point \(x\); then, by the cited result, the
morphism
\[
  G\longrightarrow M,\qquad g\longmapsto g\cdot x
\]
is smooth and therefore open, so its image is open.

Let \(U\) be the union of the equivalence classes of the closed
points of \(M\) under the action of \(G\). Then \(U\) is open and
contains every closed point of \(M\); by the Nullstellensatz,
\(U=M\). Thus \(M\) is a disjoint union of open subpreschemes, which
are necessarily also closed. Hence \(M\) is the sum prescheme of
subpreschemes \(M_i\), each of which is the \(G\)-orbit of a closed
point. Each \(M_i\) is quasi-compact, hence of finite type, and is
therefore affine by \SGARefLink{sga3:xii:5.3}{5.3}.

If \(H\) is a subgroup of multiplicative type in \(G\), it
corresponds to a \(k\)-rational point \(x\) of \(M\), and \(G/N\)
identifies with the orbit of \(x\) under \(G\), by \SGARefLink{sga3:xi:remark:5:5}{XI~5.5}~bis. Thus
\(G/N\) is affine. Since \(C\) is an open subgroup of \(N\), by
\SGARefLink{sga3:xi:proposition:5:9}{XI~5.9}, the quotient \(G/C\) is finite over \(G/N\); indeed,
\[
  G/N\simeq (G/C)/(N/C).
\]
Consequently \(G/C\) is affine as well.

This proof also gives the following result.

\paragraph{\textbf{Corollary 5.6.}}
\SGARefTarget{sga3:xii:5.6}
\emph{Under the hypotheses of \SGARefLink{sga3:xii:5.5}{5.5} on \(G\) and \(H\), the
subprescheme \(U\) of \(M\) consisting of the subgroups of
multiplicative type in \(G\) which are locally conjugate to \(H\)
is an open-and-closed subprescheme of \(M\).}

In picturesque language, every subgroup \(H'\) of multiplicative
type in \(G\) which is a limit of subgroups conjugate to \(H\) is
itself conjugate to \(H\).

\paragraph{\textbf{Remarks 5.7.}}
\SGARefTarget{sga3:xii:5.7}
Let \(S\) be a prescheme, let \(G\) be a smooth affine \(S\)-group
prescheme, let \(H\) be an \(S\)-prescheme of multiplicative type
and of finite type, and put
\[
  M=\operatorname{Hom}_{S\text{-gr}}(H,G).
\]
By \SGARefLink{sga3:xi:corollary:4:2}{XI~4.2}, \(M\) is representable, smooth over \(S\), and separated
over \(S\). One can prove for \(M\) a result entirely analogous to
\SGARefLink{sga3:xii:5.1}{5.1}, either by reducing to \SGARefLink{sga3:xii:5.1}{5.1} through an argument analogous to that
of \SGARefLink{sga3:xi:corollary:4:2}{XI~4.2}, or directly by an argument modeled on the proof of \SGARefLink{sga3:xii:5.1}{5.1}.
The corresponding variants of \SGARefLink{sga3:xii:5.3}{5.3}, \SGARefLink{sga3:xii:5.5}{5.5}, and \SGARefLink{sga3:xii:5.6}{5.6} follow; we leave
their formulation to the reader.

\section[Maximal tori and Cartan subgroups]
{Maximal tori and Cartan subgroups of not necessarily affine
algebraic groups \textup{(algebraically closed base field)}}
\SGARefTarget{sga3:xii:6}

\paragraph{\textbf{Lemma 6.1.}}
\SGARefTarget{sga3:xii:6.1}
\emph{Let \(G\) be a connected algebraic group over a field \(k\),
and let \(Z\) be an algebraic subgroup of finite index in the center
of \(G\). Then \(G/Z\) is affine.}

\emph{Proof.}
We may suppose that \(Z\) is the center of \(G\), since a scheme
finite over an affine scheme is affine. Consider the vector spaces
\[
  P^n=P^n(G)=
  \mathcal O_{G,e}/\mathfrak m_{G,e}^{\,n+1}
  \qquad (n\geq 0),
\]
where \(\mathfrak m_{G,e}\) is the maximal ideal at the identity.
The group \(G\) acts on the \(P^n(G)\) through the adjoint
representation. Let \(Z_n\) be the kernel of the corresponding
homomorphism
\[
  \operatorname{ad}_n:G\longrightarrow \operatorname{GL}(P^n).
\]
Using the connectedness of \(G\), one readily verifies that \(Z\)
is the intersection of the \(Z_n\). Since \(G\) is noetherian, this
intersection is equal to one of the \(Z_n\). Passing to the quotient,
\(\operatorname{ad}_n\) defines a monomorphism
\[
  G/Z_n\longrightarrow \operatorname{GL}(P^n),
\]
which is a closed immersion. Hence \(G/Z_n\) is affine, and so is
\(G/Z\).

\paragraph{\textbf{Lemma 6.2.}}
\SGARefTarget{sga3:xii:6.2}
\emph{Let \(G\) be a smooth algebraic group over a field \(k\), let
\(Z\) be a central algebraic subgroup, put \(G'=G/Z\), and let
\[
  u:G\longrightarrow G'
\]
be the canonical homomorphism. Let \(T\) be a subgroup of
multiplicative type in \(G\), let \(T'=u(T)\), let \(C\) be the
centralizer of \(T\) in \(G\), and let \(C'\) be the centralizer of
\(T'\) in \(G'\). Then
\[
  C'^{\,0}\subset u(C).
\]}

\emph{Proof.}
We may suppose that \(k\) is algebraically closed. Put
\[
  C_1=\bigl(u^{-1}(C'^{\,0})\bigr)_{\mathrm{red}}.
\]
It is enough to prove that \(C_1\subset C\). Indeed, \(u(C_1)\) is
connected and of finite index in \(C'^{\,0}\), hence is equal to it
set-theoretically and therefore scheme-theoretically, since \(C'\),
and consequently \(C'^{\,0}\), is smooth.

The commutator morphism
\[
  (g,t)\longmapsto gtg^{-1}t^{-1}
\]
from \(G\times G\) to \(G\) induces a morphism
\[
  \varphi:C_1\times T\longrightarrow Z_1,
  \qquad Z_1=Z_{\mathrm{red}}.
\]
Indeed, the source is reduced, and for \(g\in C_1(k)\) and
\(t\in T(k)\), the element \(gtg^{-1}t^{-1}\) belongs to \(Z(k)\)
because \(g\) centralizes \(T\) modulo \(Z\).

A calculation on \(k\)-rational points shows that
\(\varphi(g,t)\) is additive in \(g\) and in \(t\), hence
``bilinear.'' We claim that this homomorphism is identically zero;
this proves \(C_1\subset C\). By the density theorem for \(T\), we
reduce to the case in which \(T\) is finite, so that
\(n\operatorname{id}_T=0\) for some integer \(n>0\). The morphism
\(\varphi\) is then defined by a group homomorphism
\[
  C_1\longrightarrow \operatorname{Hom}_{k\text{-gr}}(T,Z_1).
\]
Since \(Z_1\) is commutative and smooth, the group on the right is
represented by an étale algebraic group over \(k\). Since \(C_1\)
is connected, every group homomorphism from \(C_1\) to that group
is zero.

\paragraph{\textbf{Corollary 6.3.}}
\SGARefTarget{sga3:xii:6.3}
\emph{Under the preceding hypotheses, suppose that \(C'\) is
connected. Then
\[
  C'=u(C),
  \qquad
  C=u^{-1}(C').
\]}

\emph{Proof.}
In this case \(C'=C'^{\,0}\). Moreover, \(C\) plainly contains
\(Z\), and hence
\[
  C=u^{-1}\bigl(u(C)\bigr).
\]

\paragraph{\textbf{Lemma 6.4.}}
\SGARefTarget{sga3:xii:6.4}
\emph{Let \(G\) be an algebraic group over a field \(k\), and let
\(Z\) be a central algebraic subgroup such that
\[
  G/Z=G'
\]
is a torus. Then \(G\) is commutative. If \(k\) is algebraically
closed, there is a torus \(T\) in \(G\) such that \(u(T)=G'\), where
\[
  u:G\longrightarrow G/Z=G'
\]
is the canonical homomorphism.}

\emph{Proof.}
We may suppose that \(k\) is algebraically closed. Consider again
the commutator morphism
\[
  G\times G\longrightarrow G,\qquad
  (g,h)\longmapsto ghg^{-1}h^{-1}.
\]
Since \(Z\) is central, this morphism factors through a morphism
\(G'\times G'\to G\). Since \(G'\) is commutative, the latter takes
its values in \(Z\), and even in \(Z_{\mathrm{red}}\), because
\(G'\times G'\) is reduced. As above, the resulting morphism
\[
  \varphi:G'\times G'\longrightarrow Z_{\mathrm{red}}
\]
is bilinear and hence zero. Thus \(G\) is commutative.

To find a torus \(T\) in \(G\) with \(u(T)=G'\), we may, after
replacing \(G\) by \(G_{\mathrm{red}}^{\,0}\), suppose that \(G\)
is smooth and connected. Replacing \(Z\) by
\(Z_{\mathrm{red}}^{\,0}\), we may also suppose that \(Z\) is
smooth and connected. By a well-known theorem of
Chevalley,\footnote{Editor's note: references should be supplied
here.} \(Z\) is an extension of an abelian variety by an affine
group. We are therefore reduced to the two cases in which \(Z\) is
affine or \(Z\) is an abelian variety. In the first case \(G\) is
affine and the result is classical \textup{(BIBLE~7, Theorem~3(a))}.

Suppose then that \(Z\) is an abelian variety. Every homomorphism
from the additive group \(\mathbb G_a\) to the torus \(G'=R\), or
to the abelian variety \(Z\), is trivial. The same is therefore true
of every homomorphism from \(\mathbb G_a\) to \(G\); thus \(G\)
contains no subgroup isomorphic to \(\mathbb G_a\). Chevalley's
structure theorem gives an exact sequence
\[
  0\longrightarrow T\longrightarrow G\longrightarrow A
  \longrightarrow 0, \tag{*}\label{sga3:xii:lemma:6-4:equation:star}
\]
where \(A\) is an abelian variety and \(T\) is a smooth connected
affine group. Since \(T\) is commutative and contains no additive
subgroup, \(T\) is a torus; it is plainly the unique maximal torus
of \(G\). It remains to prove that every epimorphism
\[
  u:G\longrightarrow R
\]
to a torus \(R\) satisfies \(u(T)=R\).

Put
\[
  \operatorname{Hom}_{\mathrm{gr}}(T,\mathbb G_m)=M,
  \qquad
  \operatorname{Hom}_{\mathrm{gr}}(R,\mathbb G_m)=P.
\]
These are finitely generated free \(\mathbb Z\)-modules, and
\[
  T=D_k(M),\qquad R=D_k(P).
\]
Let
\[
  B=\operatorname{Ext}^{1}_{k\text{-gr}}(A,\mathbb G_m);
\]
this is also the set of \(k\)-rational points of the abelian variety
dual to \(A\). We have
\[
  \operatorname{Ext}^{1}(A,T)
    =\operatorname{Hom}_{\mathrm{gr}}(M,B),
  \qquad
  \operatorname{Ext}^{1}(A,R)
    =\operatorname{Hom}_{\mathrm{gr}}(P,B). \tag{**}\label{sga3:xii:lemma:6-4:equation:double-star}
\]
In particular, the extension \(G\) of \(A\) by \(T\) is determined
by a homomorphism
\[
  \theta:M\longrightarrow B.
\]

The exact sequence \textup{(*)} also gives
\[
  0\longrightarrow \operatorname{Hom}(A,R)
  \longrightarrow \operatorname{Hom}(G,R)
  \longrightarrow \operatorname{Hom}(T,R)
  \longrightarrow \operatorname{Ext}^{1}(A,R).
\]
Moreover, \(\operatorname{Hom}(A,R)=0\), and the last map identifies
with the homomorphism
\[
  \operatorname{Hom}(P,M)\longrightarrow
  \operatorname{Hom}(P,B)
\]
induced by \(\theta:M\to B\). Setting
\[
  N=\ker\theta,
\]
we obtain canonical bijections
\[
  \operatorname{Hom}(G,R)
  \simeq \operatorname{Hom}(P,N)
  \simeq \operatorname{Hom}(S,R),
  \qquad S=D_k(N).
\]
This has the following immediate geometric description.

\paragraph{\textbf{Lemma 6.5.}}
\SGARefTarget{sga3:xii:6.5}
\emph{Let \(G\) be an extension of an abelian variety \(A\) by a
torus \(T=D_k(M)\), defined by a homomorphism
\[
  \theta:M\longrightarrow
  B=\operatorname{Ext}^{1}(A,\mathbb G_m),
\]
over an algebraically closed field \(k\). Put \(N=\ker\theta\), and
let \(S=D_k(N)\) be the corresponding torus, isomorphic to \(T/U\),
where \(U=D_k(M/N)\). Consider the extension \(H=G/U\) of \(A\) by
\(S\). This extension splits,\footnote{Editor's note: the extension
is split.} so there is a unique projection from \(H\) to \(S\), and
hence a unique homomorphism
\[
  v:G\longrightarrow S
\]
extending the canonical homomorphism \(T\to S\). For every torus
\(R\) and every homomorphism \(u:G\to R\), there is a unique
homomorphism \(u':S\to R\) such that \(u=u'v\). In particular,
\[
  \operatorname{Im}(u)=\operatorname{Im}(u')=u(T).
\]}

It follows in particular that if \(u\) is an epimorphism, then its
restriction to \(T\) is also an epimorphism. This completes the
proof of \SGARefLink{sga3:xii:6.4}{6.4}.

\paragraph{\textbf{Theorem 6.6.}}
\SGARefTarget{sga3:xii:6.6}
\emph{Let \(G\) be a smooth connected algebraic group over an
algebraically closed field \(k\).}

\begin{enumerate}
\item[\textup{\SGAInlineRefTarget{sga3:xii:theorem:6-6:item:a}\SGAInlineRefTarget{sga3:xii:6.6:item:a}(a)}] \emph{The maximal tori of \(G\) are conjugate.}

\item[\textup{\SGAInlineRefTarget{sga3:xii:6.6:item:b}\SGAInlineRefTarget{sga3:xii:theorem:6-6:item:b}(b)}] \emph{If \(T\) is a torus in \(G\), then its
centralizer is smooth and connected.}

\item[\textup{\SGAInlineRefTarget{sga3:xii:6.6:item:c}\SGAInlineRefTarget{sga3:xii:theorem:6-6:item:c}(c)}] \emph{The assignment
\[
  T\longmapsto \operatorname{Centr}_G(T)
\]
gives a bijection between the maximal tori \(T\) of \(G\) and the
Cartan subgroups of \(G\) defined in No.~1. An algebraic subgroup
\(C\) of \(G\) is a Cartan subgroup if and only if it is smooth,
nilpotent, and set-theoretically equal to its connected normalizer;
in that case it is even scheme-theoretically equal to its connected
normalizer. Moreover,
\[
  C=\operatorname{Centr}_G(T),
\]
where \(T\) is the unique maximal torus of \(C\), and
\[
  \operatorname{Norm}_G(C)=\operatorname{Norm}_G(T).
\]}

\item[\textup{\SGAInlineRefTarget{sga3:xii:6.6:item:d}(d)}] \emph{Let \(v:G\to H\) be an epimorphism of
smooth connected algebraic groups. The maximal tori, respectively
the Cartan subgroups, of \(H\) are the images of the maximal tori,
respectively the Cartan subgroups, of \(G\). If \(T\) is a maximal
torus of \(G\) and \(C\) its centralizer, then \(v(C)\) is the
centralizer of \(v(T)\).}

\item[\textup{\SGAInlineRefTarget{sga3:xii:theorem:6-6:item:e}\SGAInlineRefTarget{sga3:xii:6.6:item:e}(e)}] \emph{Under the hypotheses of \textup{\SGARefLink{sga3:xii:corollary:6-7:equation:d}{(d)}},
suppose that \(\ker v\) is a central subgroup of \(G\). Then
\[
  T\longmapsto v(T),
  \qquad
  C\longmapsto v(C)
\]
give bijections between the maximal tori, respectively the Cartan
subgroups, of \(G\) and those of \(H\). The Cartan subgroups of \(G\)
contain the center of \(G\), and hence \(\ker v\); they are precisely
the groups \(v^{-1}(C')\), where \(C'\) is a Cartan subgroup of
\(H\).}
\end{enumerate}

We also record the following immediate consequence of
\SGARefLink{sga3:xii:theorem:6-6:item:c}{Theorem~6.6(c)}.

\paragraph{\textbf{Corollary 6.7.}}
\SGARefTarget{sga3:xii:6.7}
\emph{The group \(G\) is nilpotent, that is, \(G(k)\) is nilpotent,
if and only if every torus in \(G\) is central, or equivalently if
\(G\) has a unique maximal torus. In that case this torus is the
largest subtorus of \(G\).}

\begin{proof}[Proof of \SGARefLink{sga3:xii:6.6}{Theorem 6.6}]
Let \(Z\) be a central algebraic subgroup of \(G\), put
\(G'=G/Z\), and let
\[
  u:G\longrightarrow G'
\]
be the canonical homomorphism. Then \(G'\) is smooth and connected.
If \(T'\) is a subtorus of \(G'\), it follows from \SGARefLink{sga3:xii:6.4}{6.4} that
\(u^{-1}(T')\) is commutative and that \(T'\) is the image of a
subtorus of \(u^{-1}(T')\), hence of a subtorus of \(G\). Since
\(u^{-1}(T')\) is commutative, it plainly has a largest subtorus
\(T\): the sum of two subtori is a third subtorus containing both.
Thus
\[
  u(T)=T'.
\]
It follows immediately that, for every maximal torus \(T\) of
\(G\), its image \(T'=u(T)\) is a maximal torus of \(G'\), and that
\[
  T\longmapsto u(T)
\]
is a bijection between the maximal tori of \(G\) and those of
\(G'\).

Now take
\[
  Z=\operatorname{Centr}(G)_{\mathrm{red}}^{0}.
\]
Then \(G'\) is affine by \SGARefLink{sga3:xii:6.1}{6.1}. Its maximal tori are conjugate, and
therefore so are those of \(G\), which proves \textup{\SGARefLink{sga3:xii:6.6:item:a}{(a)}}.
Moreover, \(G\) is nilpotent, respectively has a unique maximal
torus, if and only if \(G'\) has the same property. Since \(G'\) is
affine, these two conditions on \(G'\) are equivalent
\((\mathrm{BIBLE}\ 6,\ \mathrm{Th.}\ 4,\ \mathrm{Cor.}\ 2)\);
hence they are also equivalent for \(G\). If \(G\) has a unique
maximal torus \(T\), then \(T\) is invariant and therefore central;
and because every torus of \(G\) is contained in a maximal torus,
every torus is central. Conversely, if every torus is central, then
so are the maximal tori, and the conjugacy theorem
\textup{\SGARefLink{sga3:xii:6.6:item:a}{(a)}} shows that there is only one. This proves \SGARefLink{sga3:xii:6.7}{6.7}.

Let \(T\) be any torus of \(G\), let \(T'=u(T)\), and put
\[
  C'=\operatorname{Centr}_{G'}(T').
\]
The group \(C'\) is connected
\((\mathrm{BIBLE}\ 6,\ \mathrm{Th.}\ 6\textup{(a)})\).
By \SGARefLink{sga3:xii:6.3}{6.3} the centralizer \(C\) of \(T\) is
\(u^{-1}(C')\), and is therefore connected because \(Z\) is
connected. This proves \textup{\SGARefLink{sga3:xii:6.6:item:b}{(b)}}. If \(T\) is maximal, then
\(T'\) is maximal; hence \(C'\) is nilpotent, and so is \(C\),
which is a central extension of \(C'\). Moreover, \(T\) is a
maximal torus of \(C\), and by \SGARefLink{sga3:xii:6.7}{6.7} it is the unique maximal torus
of \(C\). Consequently
\[
  T\longmapsto\operatorname{Centr}_G(T)
\]
is a bijection from the maximal tori of \(G\) to the Cartan
subgroups of \(G\). We also have
\[
  \operatorname{Centr}_G(T)=C
  \subset \operatorname{Norm}_G(C)
  \subset \operatorname{Norm}_G(T).
\]
The centralizer of \(T\) has finite index in its normalizer
\((\text{cf. XI~5.9})\); the argument there remains valid without
the affineness hypothesis, because it uses only the representability
of the two functors in question, as noted in \SGARefLink{sga3:xi:unit:6:5}{XI~6.5}. Since \(C\) is
smooth and connected, it follows that
\[
  C=\operatorname{Norm}_G(C)^{0}.
\]
Furthermore, the bijection between maximal tori and Cartan
subgroups shows that the normalizers of \(C\) and \(T\) have the
same \(k\)-valued points. Since the latter normalizer is smooth,
\[
  \operatorname{Norm}_G(T)=\operatorname{Norm}_G(C).
\]

To finish the proof of \textup{\SGARefLink{sga3:xii:6.6:item:c}{(c)}}, it remains to show that a
smooth, connected, nilpotent subgroup \(C\) of \(G\) which has
finite index in its normalizer is a Cartan subgroup. Since \(Z\) is
central, the normalizer of \(C\) contains \(Z\); and since \(Z\) is
smooth and connected, it follows that \(Z\subset C\). Hence
\[
  C=u^{-1}(C'),\qquad C'=u(C),
\]
and
\[
  \operatorname{Norm}_G(C)
  =u^{-1}\bigl(\operatorname{Norm}_{G'}(C')\bigr).
\]
Thus \(C'\) is connected and nilpotent and has finite index in its
normalizer. By Theorem~1 of \(\mathrm{BIBLE}\ 7\), it is a Cartan
subgroup of \(G'\), and it follows at once that \(C\) is a Cartan
subgroup of \(G\).

We next prove \textup{\SGARefLink{sga3:xii:6.6:item:e}{(e)}}. We are in the situation considered at
the beginning of the proof, with \(\ker u=Z\) and \(G'=H\), and we
have already seen that \(T\mapsto u(T)\) is a bijection between the
maximal tori of \(G\) and those of \(G'\). Using the bijection just
proved between maximal tori and Cartan subgroups, we obtain a
bijection between Cartan subgroups \(C\) of \(G\) and Cartan
subgroups \(C'\) of \(G'\): to
\[
  C=\operatorname{Centr}_G(T)
\]
we associate
\[
  C'=\operatorname{Centr}_{G'}(T'),
  \qquad T'=u(T).
\]
Since \(C'\) is connected by \textup{\SGARefLink{sga3:xii:6.6:item:b}{(b)}}, \SGARefLink{sga3:xii:6.3}{6.3} gives
\[
  C'=u(C),\qquad C=u^{-1}(C'),
\]
which proves \textup{\SGARefLink{sga3:xii:6.6:item:e}{(e)}}.

It remains to prove \textup{\SGARefLink{sga3:xii:corollary:6-7:equation:d}{(d)}}. Let \(T\) be a maximal torus of
\(G\). We first prove that \(v(T)\) is a maximal torus of \(H\);
by the conjugacy theorem \textup{\SGARefLink{sga3:xii:6.6:item:a}{(a)}}, this will imply that all
maximal tori of \(H\) arise in this way. Let \(R\) be a torus of
\(H\) containing \(v(T)\). Replacing \(H\) by \(R\) and \(G\) by
\(v^{-1}(R)_{\mathrm{red}}^{0}\), we may suppose that \(R=H\), so
that \(H\) is a torus; it remains to prove \(v(T)=H\). Put again
\[
  Z=\operatorname{Centr}(G)_{\mathrm{red}}^{0},
  \qquad G'=G/Z,\qquad H'=H/v(Z).
\]
We have a commutative diagram with exact rows
\begin{equation}
\tag{D}\label{sga3:xii:corollary:6-7:equation:d}
\begin{tikzcd}[column sep=large,row sep=large]
e \arrow[r] &
Z \arrow[r] \arrow[d,"v''"'] &
G \arrow[r,"u"] \arrow[d,"v"'] &
G' \arrow[r] \arrow[d,"v'"'] &
e \\
e \arrow[r] &
v(Z) \arrow[r] &
H \arrow[r,"u'"'] &
H' \arrow[r] &
e .
\end{tikzcd}
\end{equation}
We already know that \(u(T)\) is a maximal torus of \(G'\).
Because \(G'\) is affine, the epimorphism
\[
  v':G'\longrightarrow H'
\]
is an epimorphism of smooth connected affine groups. Therefore
\(v'(u(T))\) is a maximal torus of \(H'\)
\((\mathrm{BIBLE}\ 7,\ \mathrm{Th.}\ 3\textup{(a)})\), and hence
equals \(H'\). Thus \(u'(v(T))=H'\), so in order to prove
\(v(T)=H\) it is enough to show that \(v(T)\supset v(Z)\). Now
\(v(Z)\) is a smooth connected subgroup of \(H\), hence a torus,
and
\[
  v'':Z\longrightarrow v(Z)
\]
is an epimorphism. By \SGARefLink{sga3:xii:6.4}{6.4},
\[
  v(Z)=v''(S)
\]
for a torus \(S\) of \(Z\). Since \(S\) is central in \(G\), it is
contained in the maximal torus \(T\), and therefore
\(v(Z)\subset v(T)\). This proves \textup{\SGARefLink{sga3:xii:6.6:item:d}{(d)}} for maximal tori.

In view of the bijection between maximal tori and Cartan subgroups,
it remains to prove that, if \(T\) is a maximal torus of \(G\) and
\(C\) its centralizer, then \(v(C)\) is the centralizer of \(v(T)\).
Return to diagram \textup{(D)}, now without assuming that \(H\) is
a torus, and put
\[
  T'=u(T),\qquad C'=u(C).
\]
By \textup{\SGARefLink{sga3:xii:6.6:item:e}{(e)}}, \(C'\) is the centralizer of the maximal torus
\(T'\). Since \(G'\), and hence \(H'\), is affine,
\(v'(C')\) is the centralizer of the maximal torus \(v'(T')\) of
\(H'\)
\((\mathrm{BIBLE}\ 7,\ \mathrm{Th.}\ 3\textup{(a)})\).
Equivalently, \(u'(v(C))\) is the centralizer of \(u'(v(T))\).
Because \(C\) contains \(Z\), the group \(v(C)\) contains \(v(Z)\).
It is therefore the inverse image under \(u'\) of \(u'(v(C))\),
that is, of the centralizer of \(u'(v(T))\). By
\textup{\SGARefLink{sga3:xii:6.6:item:e}{(e)}}, applied to \(u':H\to H'\) and the maximal torus
\(v(T)\) of \(H\), this inverse image is precisely the centralizer
of \(v(T)\). This completes the proof of \SGARefLink{sga3:xii:6.6}{6.6}.
\end{proof}
